\documentclass[french, 11pt]{article}

\input{setup.tex}

\def\chapitre{6}
\def\pagetitle{Chaînes de caractères.}

\begin{document}

\input{title.tex}

\section{Caractères.}

\begin{defi}{Caractère.}{}
    Ils sont codés sur un octet en OCaml, les normes garantissent le code ASCII.\\
    Leur type est \c{char} en C et en OCaml, le lien entre le code ASCII et le caractère peut être obtenu par conversion explicite en C, et avec \ocaml{char_of_int} ou \ocaml{int_of_char} en OCaml.\\
    Il existe des opérations entre caractères : \texttt{<, <=, >, >=, ==, !=, =, <>}.\\
    On a l'assurance que les minuscules, majuscules et chiffres soient dans l'ordre.\\
    Affichage : \texttt{"\%c"} en C, \ocaml{print_char} en OCaml. 
\end{defi}

\pagebreak

\section{Chaînes de caractères.}

Les constantes littérales seront entre guillemets.

\subsection{En C.}

\begin{defi}{}{}
    Il n'existe pas de chaîne de caractère en C, on appelle chaîne de caractères une adresse où sont stockés des caractères.\\
    Affichage : \texttt{printf("\%s", chaine);} ou bien \c{printf("chaine");}
\end{defi}

\subsection{En OCaml.}

\begin{defi}{}{}
    Le type est \ocaml{string}.\\
    Pour accéder au $i$ème caratère, on utilise \ocaml{chaine.[i]}.\\
    Les indices commencent à 0, les accès sont en temps constant. Les chaînes de caractère sont immuables.\\
    Affichage : \ocaml{print_string chaine;} ou bien \ocaml{print_endline chaine;}.\\
    Concaténation : \ocaml{chaine1 ^ chaine2}.\\
    Longueur : \ocaml{String.length chaine}.
\end{defi}

\section{Entrées et Sorties.}

\begin{defi}{Flux de données.}{}
    Un \bf{flux de données} est une transmissions ininterrompue d'une suite de données.\\
    Un processus a trois flux standards :
    \begin{itemize}
        \item \texttt{<} : entrée standard.
        \item \texttt{>} : sortie standard.
        \item \texttt{2>} : sortie d'erreur.
    \end{itemize}
\end{defi}

\subsection{En C.}

\begin{defi}{}{}
    Codes usuels :  entiers : \texttt{\%d}, flottants : \texttt{\%f}, caractères : \texttt{\%c}, pointeurs : \texttt{\%p}.\\
    Sortie standard : \c{printf(const char*, ...)}\\
    Elle bufferise par lignes.\n
    Sortie erreur : \c{fprintf(FILE*, const char*, ...);} avec \c{stderr} pour \c{FILE*}.\\
    Elle ne bufferise pas.\n
    Lire sur l'entrée standard : \c{scanf(const char*, ...);}
\end{defi}

\subsection{En OCaml.}

\begin{defi}{}{}
    Sortie standard : \texttt{print\_string}, \texttt{print\_int}, \texttt{print\_float}, ...\\
    Sortie erreur : \texttt{output\_string : out\_channel -> string -> unit} avec \texttt{stderr} pour \texttt{out\_channel}.\\
    Lire sur l'entrée standard : \texttt{read\_line}.\\
    Ces fonctions sont définies dans le module \texttt{Stdlib} ouvert par défaut.
\end{defi}

\section{Fichiers textuels.}
\subsection{En C.}

\begin{defi}{}{}
    On utilise \c{fopen(const char*, const char*);} en spécifiant le chemin du fichier et le mode d'ouverture.\\
    Lecture : \c{fscanf(FILE*, const char*, ...);}\\
    Écriture : \c{fprintf(FILE*, const char*, ...);}\\
    Fermeture : \c{fclose(FILE*);}\n
    En cas d'oubli de fermeture, on épuise les ressources du système et les écritures peuvent ne pas se faire. 
\end{defi}

\subsection{En OCaml.}

\begin{defi}{}{}
    Même principe :\\
    Ouverture : \ocaml{open_in} en lecture, \ocaml{open_out} en écriture.\\
    Lecture : \ocaml{input_line} ou bien \ocaml{input_char}.\\
    Écriture : \ocaml{output_string} ou bien \ocaml{output_char}.\\
    Fermeture : \ocaml{close_in} ou bien \ocaml{close_out}.\\
    En fin de fichier, on obtient l'exception \ocaml{End_of_file}.
\end{defi}

\section{Recherche de motifs.}

On cherche à déterminer si un motif $m$ apparaît dans un texte $T$, et où celui-ci apparaît.

\subsection{Notations.}

\begin{nota}{}{}
    Soit $T$ un texte de longueur $n$, $M$ un motif de longueur $m\leq n$, sur un même alphabet $\cursive{A}$.\\
    On utilise les notations de python sur les chaînes de caractère.\\
    On cherche $S=\{s\in\N \mid \forall i \leq m-1, ~ T[s+i]=M[i]\}$.
\end{nota}

\subsection{Algorithme naïf.}

\begin{defi}{}{}
    \begin{algorithm}[H]
        \caption{Algorithme Naïf}
        \Entree{Un texte $T$, un motif $M$}
        \Sortie{L'ensemble $S$}
        $S \gets \0$ \\
        \Pour{$s$ allant de 0 à $n-m$}{
            \Si{$M=T[s:s+m]$}{
                $S\gets S\cup\{s\}$
            }
        }
        \Retour $S$
    \end{algorithm}
    Dans le pire des cas, cet algorithme a une complexité en $\Theta(m\cdot n)$.
\end{defi}

\subsection{Algorithme de Boyer-Moore-Horspool. \texorpdfstring{$\star$}{Lg}.}
\subsubsection{Principe.}

\begin{defi}{}{}
    En comparant de droite à gauche, on peut déplacer le motif plus rapidement.\\
    On applique un prétraitement sur le motif, pour connaître la dernière occurence de chaque caractère.
\end{defi}

\subsubsection{Table de saut.}

\begin{defi}{}{}
    Indexée par l'alphabet, la case $c$ donne l'indice de l'occurence la plus à droite dans $M$ privé de sa dernière lettre du caractère $c$, sinon $-1$.\\
    Ce prétraitement s'effectue avec une complexité en $O(\max(n,m))$.
\end{defi}

\subsubsection{Algorithme.}

\begin{defi}{}{}
    \begin{algorithm}[H]
        \caption{Construction de la table de saut}
        \Entree{Un motif $M$}
        \Sortie{Table de saut sur $M$}
        $T \gets $ table de saut vide.\\
        \Pour{$c$ dans $\cursive{A}$}{
            $T[c] \gets -1$.
        }
        \Pour{$i$ allant de 0 à $m-2$}{
            $T[M[i]] \gets i$.
        }
        \Retour $T$.
    \end{algorithm}

    \begin{algorithm}[H]
        \caption{Algorithme de Boyer-Moore-Horspool}
        \Entree{Un texte $T$, un motif $M$}
        \Sortie{L'ensemble $S$}
        $S \gets \0$.\\
        $L \gets $ table de saut sur $M$.\\
        $s \gets 0$.\\
        \Tq{$s\leq n-m$}{
            $i \gets m-1$.\\
            \Tq{$i\geq 0$ et $M[i]=T[s+i]$}{
                $i \gets i-1$.
            }
            \Si{$i=-1$}{
                $S\gets S\cup\{s\}$.
            }
            $s \gets s + m - \max(1, i-L[T[s+m-1]])$.
        }
        \Retour $S$.
    \end{algorithm}
    La complexité est en $O(|T|\.|M|)$, mais en pratique, Boyer-Moore est bien plus efficace que l'algorithme naïf.
\end{defi}

\pagebreak

\subsection{Algorithme de Rabin-Karp.}

Soit $h$ sur fonction de hachage sur les chaînes de caractères, $m$ un mot de taille $n$ et $T=T_0T_1...T_{n-1}...T_l$ un texte.\\
On calcule $h(m)$ le hachage de $m$, et pour tout $i$, on le compare au hachage du mot $T_i...T_{i+n-1}$.\\
Si les hachages concordent, on teste naïvement si les mots sont égaux.\\
On peut choisir $h$ comme étant la somme (des codes) des caractères, alors :
\begin{equation*}
    h(T_{i+1}...T_{i+n}) = h(T_i...T_{i+n-1}) + h(T_{i+n}) - h(T_i).
\end{equation*}
On est alors en mesure de calculer le hachage en $O(1)$, et de le mettre à jour en $O(1)$.\\
On peut choisir d'autres $h$, mais on tient à avoir une relation de récurrence.\\
La complexité finale de l'algorithme dépendra alors de la qualité de $h$ : moins il y aura de collisions, moins il y aura de comparaisons.

\subsection{Rercherche par automate fini. (MPI)}

On construit l'automate $(\Sigma, Q, q_0, F, \d)$ reconnaissant le langage $\Sigma^* m$.\\
Lorsqu'on lit $T$, à chaque fois que l'on arrive dans un état final de l'automate, on a détecté une occurence de $m$.\\
Si on représente $\d$ sous la forme d'une matrice, on peut effectuer une transition dans l'automate en $O(1)$.\\
On peut alors lire $T$ et détecter les occurences en $O(|T|)$.\\
Pour construire l'automate, on veut, à l'état $i$, renvoyer à un état $j$ tel que $m_0...m_j$ soit le plus grand préfixe de $m_0...m_il$ qui en est aussi un suffixe. 
\begin{algorithm}
    \caption{Construction de l'automate.}
    \PourCh{$i\in\lb0,n-1\rb$}{
        \PourCh{$l\in\Sigma$}{
            \PourCh{$j\in\lb0,i+1\rb$}{
                Tester si $m_0...m_j$ est suffixe de $m_0...m_il$.
            }
            $\delta_{i,l}=j$.
        }
    }
\end{algorithm}
La construction de l'automate est donc réalisée en $O(n^3|\Sigma|)$.

\section{Compression.}
\subsection{Lempel-Ziv-Welch.}
\begin{defi}{Lempel-Ziv-Welch Compression.}{}
    \textbf{Principe:} parcours de gauche à droite du texte à compresser en maintenant un dictionnaire qui à un facteur du texte associe un code.\n
    \begin{algorithm}[H]
        \caption{LZW-Compression}
        \Entree{Un texte}
        \Sortie{Code compressé du texte}
        Créer un dictionnaire $D$.\\
        Mettre les codes des caractères dans le dictionnaire.\\
        $m \gets$ mot vide.\\
        r $\gets$ mot vide.  \emph{c'est le code qui sera renvoyé}.\\
        \Tq{il y a un caractère $x$ à lire}{
            \Si{$mx$ est une clé de $D$}{
                $m \gets mx$.
            }\Sinon{
                Ajouter $mx$ à $D$ avec la première valeur non utilisée.\\
                Ajouter le code de $m$ à la fin de $r$.\\
                $m \gets x$.
            }
        }
        Ajouter le code de $m$ à la fin de $r$.\\
        \Retour{r}.
    \end{algorithm}
\end{defi}

\vspace*{-0.8cm}

\begin{defi}{Lempel-Ziv-Welch Décompression.}{}
    \begin{algorithm}[H]
        \caption{LZW-Décompression}
        \Entree{Un code}
        \Sortie{Texte décompressé}
        Créer un dictionnaire $D$ double associant son code à chaque lettre.\\
        $c \gets$ mot vide.\\
        $r \gets $ mot vide \emph{c'est le texte qui sera renvoyé}.\\
        \Tq{il y a un code $x$ à lire}{
            \Si{$x$ est une clé de $D$}{
                $m \gets D[x]$.\\
                $m_0 \gets$ le premier caractère de $m$.\\
                \Si{$mm_0$ n'est pas une valeur de $D$}{
                    Ajouter $mm_0$ à $D$ avec la première valeur non utilisée.\\
                }
                Ajouter $m$ à la fin de $r$.\\
            }\Sinon{
                $m_0 \gets $ le premier caractère de $m$.\\
                Ajouter $mm_0$ à $D$ avec la première valeur non utilisée.\\
                Ajouter $m$ à la fin de $r$. 
            }
        }
        \Retour{r}.
    \end{algorithm}
    La décompression a seulement besoin du code, et du dictionnaire sur les caractères : on n'a pas besoin du dictionnaire complet. Elle consiste à décompresser tout en reconstruisant le dictionnaire.
\end{defi}

\subsection{Codage de Huffman.}

\begin{defi}{Huffman. $\star$}{}
    \begin{algorithm}[H]
        \caption{Codage de Huffman}
        \Entree{Un texte}
        \Sortie{Son encodage}
        Mettre les couples (caractère, fréquence) dans une chaîne de priorité.\\
        \Tq{file contient plusieurs éléments}{
            Extraire les 2 éléments de plus basses fréquences.\\
            Créer un noeud binaire dont ces éléments sont les fils.\\
            Insérer ce noeud dans la file, avec pour fréquence la somme des fréquences de ses fils.
        }
        \Retour l'élément de la file.
    \end{algorithm}
    L'algorithme de Huffman compresse lettre par lettre. À chaque lettre, on associe un code, qui est plus petit lorsque la lettre apparaît fréquemment dans le mot. Le code associé à une lettre ne peut pas être préfixe de celui d'une autre lettre.
\end{defi}

\fin



\end{document}